% Study note for RAN-31. Companion to dzw2026-vs-harmonic-null-CANONICAL.md.
% Build: tectonic -X compile design/reference/rao-score-test-STUDY-NOTE.tex
\documentclass[11pt,a4paper]{article}

\usepackage[margin=1in]{geometry}
\usepackage{amsmath,amssymb,amsthm,mathtools}
\usepackage{array}
\usepackage{booktabs}
\usepackage{enumitem}
\usepackage{microtype}
\usepackage[T1]{fontenc}
\usepackage{xcolor}
\usepackage[framemethod=default]{mdframed}
\usepackage[hidelinks]{hyperref}

\newcommand{\R}{\mathbb{R}}
\newcommand{\im}{\operatorname{im}}
\newcommand{\Hh}{\mathcal{H}}
\newcommand{\Tr}{{\mathsf T}}
\newcommand{\Pharm}{P_{\Hh}}

\theoremstyle{definition}
\newtheorem*{selfcheck}{Self-check}

\mdfdefinestyle{callout}{linewidth=0.8pt,linecolor=black!45,
  backgroundcolor=black!3,innertopmargin=8pt,innerbottommargin=8pt,
  innerleftmargin=10pt,innerrightmargin=10pt,skipabove=10pt,skipbelow=10pt}

\title{Study note: the Rao score test,\\
and why our harmonic-zero test is one}
\author{Meadowlark Bradsher}
\date{August 2026 --- companion to \texttt{dzw2026-vs-harmonic-null-CANONICAL.md}}

\begin{document}
\maketitle

\begin{mdframed}[style=callout]
\noindent\textbf{What this is.} An on-ramp. Before reading DZW \S3.3.2 you need
the Rao score test to be a thing you \emph{recognize}, not decode --- because the
whole point of the DZW reading is to confirm their machinery reduces to this
test, and you cannot check a reduction if the target is a black box. Read this
first; it makes DZW smaller.

\smallskip
\noindent\textbf{How to use it.} Work the three rungs in order, then internalize
the bridge dictionary, then read DZW. Each rung says what to \emph{extract} (one
idea), not ``read the whole book.'' There are self-checks --- if you cannot
answer them, you have not got the rung yet; that is information, not failure.

\smallskip
\noindent\textbf{Success criterion.} You own this when you can say the target
sentence (\S\ref{sec:target}) cold, in your own words, without notes. That
sentence is the whole thing compressed. Everything here is in service of being
able to say it.

\smallskip
\noindent\textbf{Honest calibration.} This is a smaller lift than DZW, and it is
genuinely accessible --- the score test is standard graduate material, written to
be taught. Doing it first is not a detour; it is the thing that makes the DZW
reading survivable. An afternoon on Rung~1, an hour on a slice of Rung~2, then
the bridge.
\end{mdframed}

\section{The one idea, before you open anything}

There are three classical ways to test a hypothesis, and you understand each by
contrast.

\begin{description}[leftmargin=1.4em,style=nextline,itemsep=0.4em]
\item[Wald] Fit the model \emph{without} the constraint, then check whether the
  unconstrained estimate is far from what the null claims.
\item[Likelihood ratio (LR)] Fit \emph{both} the constrained and unconstrained
  models, and compare how much likelihood the constraint costs.
\item[Score / Rao / Lagrange-multiplier (three names, one test)] Fit the model
  \textbf{under the null only}, then look at the log-likelihood's gradient at
  that constrained fit, \emph{in the directions the constraint forbids}. If the
  null is true, that gradient is just sampling noise. If the null is false, the
  model is straining to move into the forbidden directions, so the gradient
  there is large. Reject when it is large.
\end{description}

\noindent\textbf{Why you want the score test specifically, and not the other
two:} it is the one that \textbf{fits under the null}. That single feature is
what matches DZW, whose conditional mean is also estimated under $H_0$, and it is
what produces the ``leftover score lands in the harmonic subspace by
construction'' geometry that the rig measured at $4.1\times10^{-13}$. Wald fits
the full model and you would lose that geometry entirely. So whenever you meet
the three together, tag the score test as \emph{``the constrained-fit one.''}
That tag is why it is yours.

\section{Rung 1 --- the score test in general}

\textbf{What it says, precisely.} You have parameters split into two blocks,
$\theta=(\beta_1,\beta_2)$, and you test $H_0:\beta_2=0$, with $\beta_1$ a free
nuisance. The recipe:

\begin{enumerate}[leftmargin=1.6em,itemsep=0.45em]
\item \textbf{Fit the restricted MLE:} maximize the likelihood with $\beta_2$
  held at $0$, giving $\tilde\beta_1$. Call the whole restricted fit
  $\tilde\theta=(\tilde\beta_1,0)$.

\item \textbf{Compute the score} $U(\theta)=\partial\ell/\partial\theta$, the
  log-likelihood's gradient, at $\tilde\theta$. By stationarity of the restricted
  fit, the $\beta_1$-block of the score is $0$ at $\tilde\theta$ --- you
  maximized over $\beta_1$, so its gradient vanishes. The $\beta_2$-block $U_2$
  is generally \emph{not} zero. That is the signal.

\item \textbf{The statistic} is a quadratic form in that leftover score, with the
  Fisher information $I$ as the metric:
  \[
    T \;=\; U(\tilde\theta)^{\Tr}\,I(\tilde\theta)^{-1}\,U(\tilde\theta),
  \]
  which --- because the $\beta_1$-block is zero --- reduces to a quadratic form
  in $U_2$ alone, using the $\beta_2\beta_2$ block of $I^{-1}$, which correctly
  accounts for having estimated $\beta_1$.

\item \textbf{Under $H_0$,} $T \xrightarrow{\;d\;} \chi^2_{\dim(\beta_2)}$: the
  degrees of freedom are the number of constraints.
\end{enumerate}

\noindent\textbf{Extract this one paragraph and move on.} The score test fits
under the null, computes \emph{observed minus expected} in the forbidden
directions, and checks whether that leftover is bigger than the
Fisher-information yardstick; the df is the number of forbidden directions.

\paragraph{References.}
\begin{itemize}[leftmargin=1.4em,itemsep=0.3em]
\item Casella \& Berger, \emph{Statistical Inference} (2nd ed.), asymptotic-tests
  section ($\approx$\,\S10.3). Presents Wald, LR, and score together, which is
  the right first exposure --- you learn each by its contrast with the others.
  \emph{Primary reference for this rung.}
\item Agresti, \emph{Categorical Data Analysis}. If you want the score test in
  its logistic/binomial home turf, which is your setting, Agresti is gentler than
  the mathematical-statistics texts and the examples are pairwise and binary.
  Good complement.
\item Wasserman, \emph{All of Statistics}. Fine for the Wald test and the overall
  testing framework, but light on the score test specifically --- do not rely on
  it for this rung.
\end{itemize}

\begin{selfcheck}[Rung 1]
\begin{enumerate}[leftmargin=1.6em,itemsep=0.2em,topsep=0.3em]
\item Why is the $\beta_1$-block of the score zero at the restricted fit?
  \emph{(Because you maximized over $\beta_1$ --- stationarity.)}
\item In one sentence each, how does the score test differ from Wald and from LR?
\item Where do the degrees of freedom come from?
\end{enumerate}
\end{selfcheck}

\section{Rung 2 --- score tests in GLMs, i.e.\ ``omitted directions''}

``Omitted directions'' is not textbook vocabulary. It is the geometric name for
testing that a subset of coefficients is zero in a GLM, which classically appears
as \emph{testing a linear restriction} or \emph{comparing nested GLMs} --- the
score-test analogue of analysis of deviance. Same object, searchable under those
terms.

\paragraph{The one fact that makes your case clean --- do not skim it.}
In a GLM with the \textbf{canonical link}, the score has the tidy form
\[
  U \;=\; y-\mu \;=\; \text{observed} - \text{expected}.
\]
For the binomial with the logit link, which \emph{is} canonical, that is
\[
  U_e \;=\; w_e - k_e\,p_e,
\]
observed wins minus expected wins, per edge. The canonical link also makes the
Fisher information \textbf{diagonal} here,
\[
  I \;=\; \operatorname{diag}\!\big(k_e\,p_e(1-p_e)\big),
\]
which keeps the quadratic form simple. If you read a treatment using a
non-canonical link the score will not have this form, and you will be confused
about why yours is so clean: yours is clean \emph{precisely because} logit is
canonical for the binomial. Hold that as the reason.

\medskip
\noindent\textbf{Extract this:} ``test that these extra coefficients $\beta_2$
are zero, having fit only $\beta_1$'' is a GLM score test; under the canonical
link the leftover score is observed-minus-expected and lands in the $\beta_2$
directions; the df is $\dim(\beta_2)$. That is structurally your test --- your
``$\beta_2$ directions'' are the harmonic directions.

\paragraph{References.}
\begin{itemize}[leftmargin=1.4em,itemsep=0.3em]
\item McCullagh \& Nelder, \emph{Generalized Linear Models}. The standard. Dense;
  you need only the slice on the score test for added variables and nested
  models, and the canonical-link score form. Do not read all of it.
\item Agresti, \emph{Categorical Data Analysis} (again). Friendlier for the
  logistic case; look up ``score test'' and ``residuals for GLMs / nested
  models.''
\end{itemize}

\begin{selfcheck}[Rung 2]
\begin{enumerate}[leftmargin=1.6em,itemsep=0.2em,topsep=0.3em]
\item Why does the canonical link give score $=$ observed $-$ expected?
\item If you tested ``these 3 extra covariates are zero,'' what would the df be,
  and where would the leftover score live?
\item Why does a non-canonical link break the tidy form?
\end{enumerate}
\end{selfcheck}

\section{The bridge --- textbook to our construction}
\label{sec:bridge}

\emph{This is the core. No book contains it.}

\medskip
You will not find ``harmonic-subspace score test'' in any text, because you --- and
effectively the field --- just constructed it by \textbf{choosing the omitted
directions to be the harmonic subspace}. Here is the exact rename, line by line.
When these renames become automatic, the ``novel scary object'' collapses into
``the standard GLM score test, pointed at $\Hh$.''

\begin{center}
\newcolumntype{L}[1]{>{\raggedright\arraybackslash}p{#1}}
\begin{tabular}{@{}L{0.44\textwidth}L{0.50\textwidth}@{}}
\toprule
\textbf{Textbook GLM score test} & \textbf{Our harmonic-zero construction}\\
\midrule
natural parameter, coefficients split $(\beta_1,\beta_2)$
  & log-odds flow $\eta\in\R^{E}$, Hodge-split into
    $S=\im D_0\oplus\im D_1^{\Tr}$ (kept) and $\Hh$ = harmonic (tested-zero)\\
\addlinespace
$H_0:\beta_2=0$
  & $H_0:\eta\in S$, i.e.\ $\Pharm\,\eta=0$ --- the harmonic component of the
    true log-odds is zero (``no genuine cycle'')\\
\addlinespace
fit under $H_0$ (only $\beta_1$)
  & constrained MLE $\hat\eta=M\hat\beta$, fit inside $S$; $M$'s columns are a
    basis for $S$\\
\addlinespace
score $U=y-\mu$ (canonical link)
  & $U(\hat\eta)=w-k\,\sigma(\hat\eta)=$ observed $-$ expected wins\\
\addlinespace
$\beta_1$-block of score is $0$ at the restricted fit
  & stationarity of the constrained fit: $M^{\Tr}U(\hat\eta)=0$\\
\addlinespace
leftover score lives in the $\beta_2$ directions
  & $U(\hat\eta)\perp S$, so $U(\hat\eta)\in\Hh$ --- the harmonic subspace,
    exactly\\
\addlinespace
$\chi^2$ with $\mathrm{df}=\dim(\beta_2)$
  & $\chi^2$ with $\mathrm{df}=b_1=\dim(\Hh)$\\
\bottomrule
\end{tabular}
\end{center}

\paragraph{The two framings are the same thing.}
The textbook splits coordinates into $(\beta_1,\beta_2)$. Yours parametrizes the
kept subspace $S$ by $M$ and reads off the orthogonal complement. They coincide:
\[
  \beta_2=0
  \iff \eta\in S
  \iff \Pharm\,\eta=0,
\]
and ``leftover score in the $\beta_2$ directions'' $\iff$ $U\perp S$ $\iff$
$U\in\Hh$, because the Hodge decomposition is orthogonal, so $S^{\perp}=\Hh$. Do
not let the change of dress confuse you --- it is a reparametrization, not a
different test.

\paragraph{The by-construction harmonic landing, spelled out.}
This is the $4.1\times10^{-13}$ fact. The constrained fit satisfies
$M^{\Tr}U(\hat\eta)=0$, which is just ``the gradient vanishes in the directions
you fit over.'' The columns of $M$ span $S$, so this says $U(\hat\eta)$ is
orthogonal to \emph{all} of $S$. Orthogonal-to-$S$ is, by the orthogonal Hodge
split, exactly $\Hh$. Therefore
\[
  U(\hat\eta)\in\Hh \quad\text{exactly}
\]
--- nobody applied a harmonic projection; the stationarity condition of the
constrained fit \emph{forced} it there. The measured
\texttt{score\_off\_harmonic} $\le 4.1\times10^{-13}$ is the true value, $0$,
plus floating-point residue. That is the empirical confirmation that the geometry
did what the geometry must. ``The test lives in harmonic coordinates'' is
\textbf{literal}: the object being tested is always a vector in the
$b_1$-dimensional harmonic space, and the whole test is ``is this harmonic vector
bigger than binomial noise.''

\section{Three things not to blur while reading}

\begin{enumerate}[leftmargin=1.6em,itemsep=0.5em]
\item \textbf{Canonical link is load-bearing.} The clean observed $-$ expected
  score \emph{and} the clean orthogonal-complement landing both depend on the
  logit link being canonical for the binomial. If a source uses a different link,
  expect mess --- and know it is not your case.

\item \textbf{Wald vs.\ score.} You will meet the three tests together and it is
  easy to walk away with a mush. Tag the score test as \emph{``the
  constrained-fit one.''} That is the one your construction uses, and ``fit under
  the null'' is precisely what makes the leftover-score geometry work. Wald would
  fit the full model and destroy the by-construction landing.

\item \textbf{$I^{-1}$ does not preserve $\Hh$, so $T$ is not $\|U\|^2$
  rescaled.} The score lands in $\Hh$ exactly, but the metric does not respect
  that subspace: $I$ is diagonal in the \emph{edge} basis, not in the
  $(S,\Hh)$ split. So $T=U^{\Tr}I^{-1}U$ genuinely carries the $k_e$ and $p_e$
  dependence and cannot be read as a rescaled norm of the harmonic component.
  This matters twice --- it is where PP4's small-edge-count worry actually lives,
  and it is the step where a careless reading of DZW's studentization would
  ``simplify'' the statistic into something that is not this test.
\end{enumerate}

\section{The target sentence}
\label{sec:target}

\emph{Say this cold and you own it.}

\begin{mdframed}[style=callout]
Our test is the classical \textbf{Rao score test for omitted coefficients in a
canonical-link GLM}, where the omitted directions are chosen to be a basis for
the harmonic subspace $\Hh$. We fit the log-odds under ``no harmonic''
($\eta\in S$); the constrained fit's stationarity condition forces the leftover
score $U=\text{observed}-\text{expected}$ to be orthogonal to $S$, hence into
$\Hh$, \textbf{by construction} --- not by any projection we applied. The
statistic is the squared size of that harmonic score against the
Fisher-information metric, $\chi^2$ with $b_1$ degrees of freedom.
\end{mdframed}

\noindent When that sentence is yours without notes, RAN-31's score-test side is
owned, and DZW \S3.3.2 becomes a matter of checking whether their fold arrives at
the same statistic.

\section{Then, and only then: DZW \S3.3.2}

You are checking \textbf{one identity}: that DZW \S3.3.2, instantiated with
$H_0$ = the subspace constraint $\eta\in S$, the estimator = the $S$-constrained
MLE, and the test function $g$ = the harmonic directions, produces this same
score test. Line up four things as you read:

\begin{enumerate}[leftmargin=1.6em,itemsep=0.3em]
\item \textbf{their moment condition} --- should be ``residual $\perp g$''; with
  $g$ harmonic, it is testing harmonic-direction signal;
\item \textbf{the estimator they plug in} --- should be the constrained MLE, the
  same fit whose stationarity you now understand;
\item \textbf{the studentized statistic} --- should reduce to a quadratic form in
  the held-out harmonic score;
\item \textbf{the df} --- should come out $b_1$.
\end{enumerate}

If those align, DZW in the fixed-graph case \emph{is} the test above ---
``checking a thing against itself'' --- and you can say so in your own words:
both roads fit under ``no harmonic'' and test the leftover harmonic signal
against noise; the score test does it in-sample via stationarity, DZW does it
cross-fit via folds, and in the fixed-graph case they are the same statistic.

\paragraph{Where to stop --- this is permission, not a limit.}
There is a further piece in DZW \S4: the debiasing correction, the
geometric-series $(I+B)^{-1}$ factor. That is \textbf{not today's job}. It is for
the case where the harmonic loops $g$ are chosen \emph{from the data} (adaptive,
post-selection), where the score test is invalid and the fold is what restores
validity. That is the genuinely hard part, it is the reason DZW matters beyond
the fixed case, and it is the right question to bring the specialist.

So if you follow \S3.3.2 to the fixed-case identity, then hit the $(I+B)^{-1}$
debiasing and lose the thread --- \textbf{that is a finished, successful
reading.} You will have confirmed what RAN-31 needs for the fixed-graph claim and
located the exact boundary:

\begin{quote}
\emph{``I follow it up to the debiasing correction for the selected case, and
that's where I need you.''}
\end{quote}

\noindent That sentence is the sharpest thing you could walk into the seminar
with.

\section{Consolidated self-check}

\emph{The ``advanced undergrad asking why'' pass. If you can answer these without
notes, you are ready for DZW.}

\begin{enumerate}[leftmargin=1.6em,itemsep=0.35em]
\item \textbf{Why the score test and not Wald or LR?}
  \emph{(Fits under the null, so it matches DZW and gives the by-construction
  geometry.)}
\item \textbf{Why is the leftover score orthogonal to $S$?}
  \emph{(Constrained-fit stationarity: $M^{\Tr}U=0$.)}
\item \textbf{Why is orthogonal-to-$S$ the same as harmonic?}
  \emph{(The Hodge split is orthogonal, so $S^{\perp}=\Hh$.)}
\item \textbf{Why is the score observed $-$ expected here?}
  \emph{(Canonical, i.e.\ logit, link.)}
\item \textbf{Why $\mathrm{df}=b_1$?}
  \emph{($b_1=\dim\Hh$ = number of omitted directions = codimension of the
  constraint.)}
\item \textbf{What is \texttt{score\_off\_harmonic} $\le 4.1\times10^{-13}$
  confirming?} \emph{(That $U(\hat\eta)\in\Hh$ is exact; the tiny number is
  floating-point residue, not a modelling choice.)}
\item \textbf{Why is $T$ not just the squared norm of the harmonic score?}
  \emph{($I^{-1}$ does not preserve $\Hh$; the metric carries $k_e$ and $p_e$.)}
\item \textbf{What are you checking in DZW \S3.3.2, and where do you stop?}
  \emph{(Whether their fold gives the same statistic; stop at the \S4
  selected-case debiasing.)}
\end{enumerate}

\end{document}
